% ==========================================================================
% Supervisor Meeting Agenda - 2 March 2026 (30 minutes)
%
% Compile with: pdflatex agenda.tex
% ==========================================================================

\documentclass[11pt,a4paper]{article}
\usepackage[utf8]{inputenc}
\usepackage[margin=2.5cm]{geometry}
\usepackage{enumitem}
\usepackage{hyperref}
\usepackage{booktabs}
\usepackage{tabularx}
\usepackage{xcolor}

\definecolor{done}{RGB}{34,139,34}
\definecolor{todo}{RGB}{200,50,50}
\definecolor{question}{RGB}{0,90,180}

\newcommand{\done}[1]{\textcolor{done}{\textbf{DONE:}} #1}
\newcommand{\todo}[1]{\textcolor{todo}{\textbf{TODO:}} #1}
\newcommand{\question}[1]{\textcolor{question}{\textbf{Q:}} #1}

\title{Supervisor Meeting Agenda\\
\large 2 March 2026 --- 30 minutes}
\author{Antonio Galdes}
\date{}

\begin{document}
\maketitle

% ------------------------------------------------------------------
\section*{1. Progress Update (5 min)}
% ------------------------------------------------------------------

\begin{itemize}[leftmargin=*]
  \item \done{Task 1 --- Unit tests and CI pipeline (all passing, GitHub Actions green)}
  \item \done{Task 2 --- Linting, formatting, type checking}
  \item \done{Task 3 --- Evidential actor integrated into SB3
    (\texttt{EvidentialPPO}, \texttt{EvidentialActorCriticPolicy}, NIG distributions)}
  \item \done{Task 4 --- Baseline configs (4 YAML overrides) + orchestration script
    (\texttt{run\_experiment.py}) for ablation study}
  \item \done{PPO hyperparameter review --- calibrated against literature
    (Andrychowicz et al.\ 2021, CaRL 2025, EPPO 2024)}
  \item 154 unit tests passing, \texttt{make verify} green (lint + format + typecheck + tests)
  \item All code runs on CPU --- zero GPU time used so far
\end{itemize}

% ------------------------------------------------------------------
\section*{2. Decisions Needed (15 min)}
% ------------------------------------------------------------------

\subsection*{2.1 EKF Sensor Inputs (Task 7)}

The agent's observation comes from the EKF. The EKF can fuse any subset of the sensors below. Need to decide which combination to use --- affects drift behaviour, covariance characteristics, and what is feasible on the real-world deployment vehicles.

\begin{table}[h]
\centering
\small
\begin{tabularx}{\textwidth}{@{}lllX@{}}
\toprule
\textbf{Sensor} & \textbf{Provides} & \textbf{Drift?} & \textbf{Notes} \\
\midrule
Wheel odometry & Velocity, displacement & Accumulates & Always available on real vehicles \\
\addlinespace
IMU & Angular vel., linear accel. & Gyro bias drift & Cheap, always available \\
\addlinespace
GNSS (GPS/Galileo) & Absolute position & No drift (noisy) & Unreliable near buildings, unusable indoors \\
\addlinespace
Camera (visual odom.) & Relative motion from features & Accumulates & Requires feature extraction; lighting-dependent \\
\addlinespace
Ray-cast (sim. LiDAR) & Distance to surfaces & No drift & Requires scan-matching (e.g.\ Cartographer) \\
\addlinespace
LiDAR & Point cloud & No drift & Dense; requires real hardware for deployment \\
\bottomrule
\end{tabularx}
\end{table}

\question{Which subset of sensors should the EKF fuse? What is available on the deployment vehicles? What produces the most interesting covariance signal for the thesis without excessive implementation scope?}

\bigskip
\subsection*{2.2 CARLA Map and Parking Scenarios (Task 5)}

Current plan: use CARLA Town01 parking lot with 3 scenario types:

\begin{enumerate}[leftmargin=*]
  \item Perpendicular parking (90\textdegree{} spots, 2.5m $\times$ 5m)
  \item Parallel parking (street-side, 2.5m $\times$ 6m)
  \item Angled parking (45\textdegree{} spots, 2.5m $\times$ 5.5m)
\end{enumerate}

\question{Should I use a different or larger map? Should I start with one scenario type or implement all three from the beginning?}

\question{Forward-entry only (current action space has no reverse gear) --- is this acceptable?}

\subsection*{2.3 Obstacle Perception}

Without perception, the agent memorises the training layout and cannot detect novel obstacles. Some form of spatial awareness is needed so the agent can generalise across layouts.

\begin{table}[h]
\centering
\small
\begin{tabularx}{\textwidth}{@{}llXl@{}}
\toprule
\textbf{Option} & \textbf{Observation} & \textbf{Weather affects it?} & \textbf{Scope} \\
\midrule
Ray-cast / ultrasonic & 4--8 distance floats & No & Minimal \\
LiDAR (sector-reduced) & 8--16 min-distance floats & Yes (fog, rain) & Moderate \\
LiDAR (full point cloud) & Thousands of 3D points & Yes & Significant \\
Camera (RGB) & Images (640$\times$480$\times$3) & Yes (fog, glare) & Massive \\
\bottomrule
\end{tabularx}
\end{table}

\question{Which perception modality? Weather-immune sensors (ultrasonic/ray-cast) isolate the localisation uncertainty variable cleanly. Weather-dependent sensors (LiDAR/camera) introduce a confounding variable but are more realistic.}

\subsection*{2.4 Teleoperation Handoff Evaluation (Task 11)}

The thesis claims the agent ``drives more carefully when uncertain''. One safety mechanism is handing off to a human teleoperator when evidential uncertainty exceeds a threshold.

\question{How should I evaluate handoff outcomes? Switching to manual control is not necessarily bad --- the model might also handle high uncertainty by driving cautiously on its own. What should the comparison look like?}

\subsection*{2.5 Compute Resources and Seeds}

Full ablation study: 4 baselines $\times$ N scenarios $\times$ N seeds.

\textbf{Per-seed time estimate} (assumptions below):

\begin{table}[h]
\centering
\small
\begin{tabular}{@{}ll@{}}
\toprule
\textbf{Parameter} & \textbf{Value} \\
\midrule
Timesteps per seed & 1M \\
Episode length & 100--300 steps (parking is short; max 500) \\
CARLA sync tick rate & 50--200 steps/sec (bottleneck is CARLA, not policy) \\
Scene & Single parking lot, low-poly, vector obs only (no camera rendering) \\
Weather + NPCs & Randomised per episode; 0--5 vehicles, 0--3 pedestrians \\
\midrule
\textbf{Est.\ per seed} & \textbf{3--6 hours} \\
\bottomrule
\end{tabular}
\end{table}

\textbf{Parallelisation strategy --- two levels available:}

\begin{enumerate}[leftmargin=*]
  \item \textbf{Run-level parallelism (implemented):} Multiple independent training runs execute simultaneously, each with its own CARLA instance on a separate port. Seeds are embarrassingly parallel --- no coordination needed. Controlled via \texttt{parallel\_workers} in config or \texttt{--workers} CLI flag.

  \textit{Pros:} Simple, already implemented, no impact on training dynamics.\\
  \textit{Cons:} Each worker needs its own CARLA instance ($\sim$2--4 GB VRAM each).

  \item \textbf{Env-level parallelism (not implemented, optional):} Multiple CARLA environments feed into a \emph{single} PPO model within one training run (SB3 \texttt{SubprocVecEnv}). The rollout buffer fills faster because data is collected from multiple envs simultaneously.

  \textit{Pros:} Each individual run finishes faster.\\
  \textit{Cons:} Doubles CARLA instances needed per run; changes training dynamics (PPO hyperparameters like \texttt{n\_steps} and \texttt{batch\_size} were calibrated for a single env); more complex to implement.
\end{enumerate}

\textbf{Current state:} Run-level parallelism is implemented and ready. Env-level parallelism is not needed unless individual runs prove too slow.

\begin{table}[h]
\centering
\small
\begin{tabular}{@{}lllll@{}}
\toprule
\textbf{Seeds} & \textbf{Total runs} & \textbf{1 worker} & \textbf{2 workers} & \textbf{3 workers} \\
\midrule
5 & 20 & 3--5 days & 1.5--2.5 days & 1--2 days \\
10 & 40 & 5--10 days & 2.5--5 days & 2--3 days \\
15 & 60 & 8--15 days & 4--7.5 days & 3--5 days \\
\bottomrule
\end{tabular}
\end{table}

\textit{Note: parking-only scenes with vector observations (no camera rendering) are lightweight. Each CARLA instance needs $\sim$2--4 GB VRAM, so a 24 GB GPU can run 2--3 instances in parallel. Actual throughput depends on GPU model.}

\question{What GPU hardware is available to me and what do I need to do to get access? How many parallel CARLA instances can the hardware support? Is 5 seeds acceptable if compute is limited?}

% ------------------------------------------------------------------
\section*{3. Timeline and Next Steps (5 min)}
% ------------------------------------------------------------------

\begin{table}[h]
\centering
\small
\begin{tabular}{@{}lll@{}}
\toprule
\textbf{Task} & \textbf{Est.\ Time} & \textbf{Blocked on} \\
\midrule
Task 6: Docker stack verification & 1--2 days & GPU machine \\
Task 5: CARLA parking environment & 1 week & \textbf{Map decision (today)} \\
Task 7: EKF filtered observation & 3--5 days & \textbf{Sensor decision (today)} \\
Task 8: End-to-end training & 2--3 days & Tasks 5, 6, 7 \\
Task 9: Reward function tuning & 1--2 weeks & Task 8 \\
Training runs (4 baselines $\times$ 10 seeds) & 2--4 weeks & Task 9 \\
Tasks 10--14: Evaluation + analysis & 2 weeks & Trained models \\
Dissertation writing & Ongoing & Can start now \\
\bottomrule
\end{tabular}
\end{table}

\question{Does this timeline align with the submission deadline? Any priorities to adjust?}

% ------------------------------------------------------------------
\section*{4. Dissertation Writing (5 min)}
% ------------------------------------------------------------------

Starting to write chapters that do not depend on results:

\begin{itemize}[leftmargin=*]
  \item Introduction --- problem statement, contributions
  \item Background --- evidential deep learning, EKF localisation, PPO, uncertainty in RL
  \item Methodology --- 2$\times$2 ablation design, architecture, observation space
  \item Experimental setup --- CARLA environment, Docker stack, evaluation conditions
\end{itemize}

\question{Any guidance on dissertation structure, expected length, or formatting requirements?}

\end{document}
